\documentclass[10pt,aspectratio=169]{beamer}
% Preámbulo modular para presentaciones Beamer (Modelación Estadística)
% Diseño y paleta institucional del Tecnológico de Monterrey

\documentclass[aspectratio=169,xcolor={dvipsnames,table}]{beamer}

\usepackage[utf8]{inputenc}
\usepackage[spanish,mexico]{babel}
\usepackage{amsmath,amssymb,amsfonts,mathtools}
\usepackage{graphicx}
\usepackage{listings}
\usepackage{booktabs}
\usepackage{tikz}
\usetikzlibrary{matrix,arrows,decorations.pathmorphing,shapes.geometric,positioning}

% Paleta Institucional Tecnológico de Monterrey
\definecolor{TecAzul}{HTML}{0039A6}       % Azul TEC: color primario institucional
\definecolor{TecAzulOscuro}{HTML}{1A2E51} % Azul marino oscuro
\definecolor{TecRojo}{HTML}{EC2661}       % Rojo de acento (paleta miTec)
\definecolor{TecGris}{HTML}{646464}       % Gris neutro para comentarios y subtítulos
\definecolor{TecFondoBloque}{HTML}{F4F6F9}% Fondo suave para bloques

% Configuración de Tema Beamer
\usetheme{Boadilla}
\usecolortheme[named=TecAzulOscuro]{structure}
\setbeamercolor{alerted text}{fg=TecRojo}
\setbeamercolor{palette primary}{bg=TecAzulOscuro,fg=white}
\setbeamercolor{palette secondary}{bg=TecAzul,fg=white}
\setbeamercolor{palette tertiary}{bg=TecAzulOscuro!80!black,fg=white}
\setbeamercolor{title}{fg=TecAzulOscuro,bg=TecFondoBloque}
\setbeamercolor{frametitle}{fg=TecAzulOscuro,bg=TecFondoBloque}

% Bloques personalizados elegantes
\setbeamercolor{block title}{bg=TecAzulOscuro,fg=white}
\setbeamercolor{block body}{bg=TecFondoBloque,fg=black}
\setbeamercolor{block title alerted}{bg=TecRojo,fg=white}
\setbeamercolor{block body alerted}{bg=TecRojo!8,fg=black}
\setbeamercolor{block title example}{bg=TecAzul,fg=white}
\setbeamercolor{block body example}{bg=TecAzul!8,fg=black}

% Redondear bordes y añadir sombra sutil
\setbeamertemplate{blocks}[rounded][shadow=true]
\setbeamertemplate{navigation symbols}{}

% Pie de página limpio con número de diapositiva
\setbeamertemplate{footline}{
  \leavevmode%
  \hbox{%
  \begin{beamercolorbox}[wd=.7\paperwidth,ht=2.25ex,dp=1ex,left]{author in head/foot}%
    \hspace*{2ex}\insertshortauthor~~(\insertshortinstitute) \hfill \insertshorttitle
  \end{beamercolorbox}%
  \begin{beamercolorbox}[wd=.3\paperwidth,ht=2.25ex,dp=1ex,right]{date in head/foot}%
    \insertshortdate{}\hspace*{2em}
    \insertframenumber{} / \inserttotalframenumber\hspace*{2ex}
  \end{beamercolorbox}}%
  \vskip0pt%
}

% Configuración de Listings de Python para visualización pedagógica de código
\lstset{
    language=Python,
    basicstyle=\ttfamily\footnotesize,
    keywordstyle=\color{TecAzulOscuro}\bfseries,
    stringstyle=\color{TecRojo},
    commentstyle=\color{TecGris}\itshape,
    showstringspaces=false,
    breaklines=true,
    frame=lines,
    rulecolor=\color{TecAzulOscuro!30},
    backgroundcolor=\color{TecFondoBloque},
    aboveskip=0.5em,
    belowskip=0.5em,
    literate={á}{{\'a}}1 {é}{{\'e}}1 {í}{{\'i}}1 {ó}{{\'o}}1 {ú}{{\'u}}1
             {Á}{{\'A}}1 {É}{{\'E}}1 {Í}{{\'I}}1 {Ó}{{\'O}}1 {Ú}{{\'U}}1
             {ñ}{{\~n}}1 {Ñ}{{\~N}}1 {¿}{{?`}}1 {¡}{{!`}}1
}

% Metadatos institucionales por defecto
\title[Modelación Estadística]{Modelación Estadística para Ciencia de Datos e Ingeniería}
\author[J. Castillo Colmenares]{Juliho David Castillo Colmenares}
\institute[Tec de Monterrey]{Tecnológico de Monterrey \\ Escuela de Ingeniería y Ciencias}
\date{\today}
% Comandos y macros estadísticas compartidas para las presentaciones Beamer (Metropolis)

% Conjuntos numéricos básicos
\newcommand{\R}{\mathbb{R}}
\newcommand{\N}{\mathbb{N}}
\newcommand{\Q}{\mathbb{Q}}
\newcommand{\Z}{\mathbb{Z}}
\newcommand{\set}[1]{\left\{ #1 \right\}}
\newcommand{\abs}[1]{\left|#1\right|}
\newcommand{\norm}[1]{\left\| #1 \right\|}

% Comandos estadísticos y probabilísticos del curso
\newcommand{\Var}{\operatorname{Var}}
\newcommand{\cov}{\operatorname{Cov}}
\newcommand{\corr}{\rho}
\newcommand{\comb}[2]{\begin{pmatrix} #1 \\ #2 \end{pmatrix}}
\newcommand{\s}{\sigma}
\newcommand{\particion}{\mathfrak{P}}
\newcommand{\card}[1]{\#\left( #1 \right)}
\newcommand{\seq}[3]{\set{#1}_{#2}^{#3}}
\newcommand{\E}{\mathbb{E}}
\newcommand{\Prob}{\mathbb{P}}

% Entornos pedagógicos y cajas en español académico natural
\newenvironment{discusion}{%
  \begin{alertblock}{Pregunta de Discusión}%
}{%
  \end{alertblock}%
}

\newenvironment{reflexion}{%
  \begin{alertblock}{Reflexión para el Aula}%
}{%
  \end{alertblock}%
}

\newenvironment{paradoja}{%
  \begin{alertblock}{Paradoja de Exploración}%
}{%
  \end{alertblock}%
}

\newenvironment{motivacion}{%
  \begin{exampleblock}{Motivación: Ciencia de Datos en la Práctica}%
}{%
  \end{exampleblock}%
}

\newenvironment{retoclase}{%
  \begin{block}{Reto Rápido en Equipo}%
}{%
  \end{block}%
}

\title[Bloques y cuadrados latinos]{Sección 08.06: Bloques, cuadrados latinos y grecolatinos}\subtitle{Control de fuentes de variabilidad}\author[J. Castillo Colmenares]{Juliho Castillo Colmenares}\institute{Tecnológico de Monterrey}\date{}
\begin{document}
\begin{frame}[plain]\titlepage\end{frame}
\begin{frame}{Hoja de ruta}\small\begin{enumerate}\item DBCA y modelo aditivo.\item Cuadrado latino.\item Cuadrado grecolatino.\end{enumerate}\end{frame}
\begin{frame}{Fuente y alcance}\small La sección introduce diseños que controlan una, dos o tres fuentes conocidas de heterogeneidad además del tratamiento.\begin{block}{Pregunta guía}¿Cómo reducir ruido sin multiplicar innecesariamente las corridas?\end{block}\end{frame}
\begin{frame}{Motivación del bloqueo}\small Si las unidades difieren de forma conocida, un diseño completamente aleatorizado mezcla esa variación con el error. Los bloques la separan.\end{frame}
\begin{frame}{DBCA}\small En un Diseño en Bloques Completos al Azar hay $b$ bloques homogéneos y $k$ tratamientos; cada tratamiento aparece una vez por bloque, $N=kb$.\end{frame}
\begin{frame}{Modelo aditivo DBCA}\small\[Y_{ij}=\mu+\tau_i+\beta_j+\epsilon_{ij}.\]\begin{itemize}\item $\tau_i$: tratamiento.\item $\beta_j$: bloque.\item $\epsilon_{ij}$: error.\end{itemize}\end{frame}
\begin{frame}{Partición DBCA}\small\[\mathrm{SCT}=\mathrm{SCTR}+\mathrm{SCB}+\mathrm{SCE}.\]La variación del bloque se extrae antes de evaluar tratamientos.\end{frame}
\begin{frame}{Cuadrado latino}\small Un cuadrado latino de orden $n$ es una matriz $n\times n$ donde cada tratamiento aparece exactamente una vez en cada fila y columna.\end{frame}
\begin{frame}{Dos fuentes de bloqueo}\small Filas y columnas representan dos factores de control cruzados. El diseño usa $n^2$ corridas en lugar de las $n^3$ de un factorial completo de tres factores.\end{frame}
\begin{frame}{Modelo latino}\small\[Y_{ijk}=\mu+\rho_i+\gamma_j+\tau_k+\epsilon_{ijk}.\]\end{frame}
\begin{frame}{Partición latina}\small\[\mathrm{SCT}=\mathrm{SC}_{\mathrm{filas}}+\mathrm{SC}_{\mathrm{columnas}}+\mathrm{SCTR}+\mathrm{SCE}.\]\end{frame}
\begin{frame}{Cuadrado grecolatino}\small Superpone dos cuadrados latinos ortogonales: el arreglo latino representa el tratamiento principal y las letras griegas una cuarta fuente controlada.\end{frame}
\begin{frame}{Condición de existencia}\small No existen cuadrados grecolatinos ortogonales de orden $6$; para otros órdenes $n\geq3$ sí se dispone de construcciones.\end{frame}
\begin{frame}{Modelo grecolatino}\small\[Y_{ijkl}=\mu+\rho_i+\gamma_j+\tau_k+\delta_l+\epsilon_{ijkl}.\]\end{frame}
\begin{frame}{Partición grecolatina}\small\[\mathrm{SCT}=\mathrm{SC}_{\mathrm{filas}}+\mathrm{SC}_{\mathrm{columnas}}+\mathrm{SCTR}+\mathrm{SC}_{\mathrm{griego}}+\mathrm{SCE}.\]\end{frame}
\begin{frame}{Aplicación: DBCA}\small Para comparar tratamientos en servidores de distintas generaciones, use cada generación como bloque y aleatorice los tratamientos dentro de ella.\end{frame}
\begin{frame}{Aplicación: cuadrado latino}\small Si turno y máquina son dos fuentes de ruido cruzadas, filas y columnas pueden controlarlas mientras cada tratamiento aparece una vez en cada dirección.\end{frame}
\begin{frame}{Aplicación: lectura de medias}\small Las medias de tratamiento se comparan después de sustraer las contribuciones sistemáticas de bloques, filas o columnas.\end{frame}
\begin{frame}{Interpretación}\small Un bloque eficaz reduce CME y puede aumentar la potencia del estadístico $F$ para tratamientos.\end{frame}
\begin{frame}{Comunicación}\small Reporte estructura, restricciones, sumas de cuadrados, grados de libertad y la hipótesis de tratamientos; los bloques son controles, no necesariamente objetivos.\end{frame}
\begin{frame}{Síntesis}\small\begin{itemize}\item DBCA controla una fuente.\item Latino controla dos.\item Grecolatino controla tres.\item El objetivo es reducir el error residual.\end{itemize}\end{frame}
\begin{frame}{Conexión con el capítulo}\small Estos diseños extienden ANOVA y preparan el análisis de diseños factoriales con más de un factor de tratamiento.\end{frame}
\end{document}
